% Generálva: html/build_tex.js (o08.js). Kézzel ne szerkeszd, a forrás a html/data mappa.
\setcounter{chapter}{7}
\chapter{\texorpdfstring{Műsorszórás és digitális TV}{Műsorszórás és digitális TV}}\label{f8}
Szuperheterodin vevő és tükörfrekvencia, valamint a DVB-T multiplex nettó bitsebessége.

\section{Alapok, egyszerűen}
\subsection{\texorpdfstring{A1. Szuperheterodin vevő}{A1. Szuperheterodin vevő}}\label{b8-szuper}
A vevő a venni kívánt $f_v$ frekvenciájú jelet egy helyi oszcillátor ($f_\text{osc}$) jelével keveri (szorozza). A szorzat a különbségi és az összegfrekvencián jelenik meg. A különbségi frekvencia az állandó \textbf{középfrekvencia} ($f_\mathrm{KF}$), ezt erősíti és szűri a vevő, az összeget eldobja:

\begin{center}\adjustbox{max width=\linewidth}{$\displaystyle |f_v-f_\text{osc}|=f_\mathrm{KF}$}\end{center}

\textbf{Felső keverés:} $f_\text{osc}=f_v+f_\mathrm{KF}$. \textbf{Alsó keverés:} $f_\text{osc}=f_v-f_\mathrm{KF}$.

A \textbf{tükörfrekvencia} ($f_t$) az a másik frekvencia, amelyik ugyanazzal az oszcillátorral szintén $f_\mathrm{KF}$-et ad: az oszcillátor túloldalán, tőle $f_\mathrm{KF}$-re. Felső keverésnél $f_t=f_\text{osc}+f_\mathrm{KF}=f_v+2f_\mathrm{KF}$, alsónál $f_t=|f_\text{osc}-f_\mathrm{KF}|=|f_v-2f_\mathrm{KF}|$. A tükröt a keverő \textbf{előtt} kell kiszűrni, utána már nem különböztethető meg a hasznos jeltől.

\subsection{\texorpdfstring{A2. DVB-T (OFDM)}{A2. DVB-T (OFDM)}}\label{b8-dvb}
Sok keskeny vivő egymás mellett, mindegyik egy-egy QAM-szimbólumot visz. $T_0$ a hasznos szimbólumidő, a vivőköz $1/T_0$. Minden szimbólum elé védőintervallum kerül ($GI$, a $T_0$ hányadában), így egy OFDM-szimbólum $(1+GI)T_0$ ideig tart. A sávszélesség kb. $(N_\text{vivő}+1)/T_0$. Csak a hasznos vivők visznek adatot:

\begin{center}\adjustbox{max width=\linewidth}{$\displaystyle R_s=\frac{N_\text{hasznos}}{(1+GI)T_0}\qquad v_\text{nettó}=R_s\cdot\log_2M\cdot\text{FEC}\cdot\frac{188}{204}$}\end{center}

A FEC a konvolúciós kód aránya, a $188/204$ az RS(204,188) kódé (204 bájtból 188 a hasznos).

% ------------------------------------------------------------------
\section{\texorpdfstring{8.1. feladat: Felső keverés: 97,6 és 107 MHz}{8.1. feladat: Felső keverés: 97,6 és 107 MHz}}\label{o8_01}
\setcounter{lepes}{0}
\begin{feladat}{8.1. feladat}
KF 10,7 MHz, felső keverés. Oszcillátor és tükör a 97,6 és a 107 MHz-es adóhoz?

{\footnotesize\sffamily Forrás: 5-ösöm 24. oldal $\cdot$ 5o\_24\_2, 5o\_24\_3. Címkék: vizsga}
\end{feladat}
\begin{keresunk}
\begin{itemize}
\item $f_\text{osc}$: az oszcillátor frekvenciája
\item $f_t$: tükörfrekvencia
\end{itemize}
\end{keresunk}
\begin{jeloles}
\begin{jeltabla}
$f_\mathrm{KF}$ & középfrekvencia & \\
$\ca{f_v}$ & a venni kívánt adó frekvenciája & $97{,}6$\\
$f_\text{osc}$ & az oszcillátor frekvenciája & \emph{keresett}\\
$f_t$ & tükörfrekvencia & \emph{keresett}\\
\end{jeltabla}
\par\smallskip\textbf{Adott még:}
\begin{itemize}
\item felső keverés: az oszcillátor a vett frekvencia fölött
\end{itemize}
\end{jeloles}

\lepes{97,6 MHz}
Felső keverésnél az oszcillátor $f_\mathrm{KF}$-fel a vett frekvencia felett van, a tükör még egyszer $f_\mathrm{KF}$-fel feljebb (A1).
\behely{\bh{\ca{f_v}}{\ca{97{,}6}}}

\begin{center}\adjustbox{max width=\linewidth}{$\displaystyle f_\text{osc}={\ca{97{,}6}}+10{,}7=108{,}3\ \text{MHz}$}\end{center}
\jegyzet{behelyettesítés}{}

\begin{center}\adjustbox{max width=\linewidth}{$\displaystyle f_t=108{,}3+10{,}7=119\ \text{MHz}$}\end{center}
\jegyzet{behelyettesítés}{ellenőrzés: $119-108{,}3=10{,}7$}

\lepes{107 MHz}
Ugyanígy.

\begin{center}\adjustbox{max width=\linewidth}{$\displaystyle f_\text{osc}=107+10{,}7=117{,}7\ \text{MHz}\qquad f_t=117{,}7+10{,}7=128{,}4\ \text{MHz}$}\end{center}
\jegyzet{behelyettesítés}{}
\begin{eredmeny}

\begin{center}\adjustbox{max width=\linewidth}{$\displaystyle 97{,}6\ \text{MHz: }f_\text{osc}=108{,}3,\ f_t=119\ \text{MHz}\qquad107\ \text{MHz: }f_\text{osc}=117{,}7,\ f_t=128{,}4\ \text{MHz}$}\end{center}
\end{eredmeny}
\begin{tipp}[Ellenőrzés]
\begin{itemize}
\item Mindkét tükör a CCIR FM-sáv (87,5 és 108 MHz) fölött van, ezért egy sávszűrő kiszűri. Ezért választják a 10,7 MHz-et: $2f_\mathrm{KF}=21{,}4$ MHz nagyobb a sáv szélességénél (20,5 MHz).
\end{itemize}
\end{tipp}


% ------------------------------------------------------------------
\section{\texorpdfstring{8.2. feladat: 24\ldots{}36 MHz, KF 18 MHz}{8.2. feladat: 24…36 MHz, KF 18 MHz}}\label{o8_02}
\setcounter{lepes}{0}
\begin{feladat}{8.2. feladat}
A vételi sáv 24 és 36 MHz között van, a KF 18 MHz. a) Az oszcillátor tartománya felső és alsó keverésnél, melyik jobb? b) Alsó keverésnél a 25,5 MHz-hez? c) Tükörfrekvencia, mi vele a gond?

{\footnotesize\sffamily Forrás: 5-ösöm 24. oldal $\cdot$ 5o\_24\_4. Címkék: vizsga, javítás}
\end{feladat}
\begin{keresunk}
\begin{itemize}
\item a) $f_\text{osc}$ (az oszcillátor frekvenciája) tartománya mindkét esetben
\item b) $f_\text{osc}$ (az oszcillátor frekvenciája) a 25,5 MHz-hez, alsó keverés
\item c) $f_t$ (tükörfrekvencia) és a gond vele
\end{itemize}
\end{keresunk}
\begin{jeloles}
\begin{jeltabla}
$\ca{f_v}$ & a venni kívánt adó frekvenciája & $24$\,MHz\\
$f_\mathrm{KF}$ & középfrekvencia & \\
$f_\text{osc}$ & az oszcillátor frekvenciája & \emph{keresett}\\
$f_t$ & tükörfrekvencia & \emph{keresett}\\
$f$ & a keresett zavaró (tükör) frekvencia & \\
\end{jeltabla}
\end{jeloles}

\lepes{a) Hangolási tartomány}
A sáv két szélére kiszámoljuk az oszcillátort (A1). Egy oszcillátort annál könnyebb megépíteni, minél kisebb a legnagyobb és a legkisebb frekvencia aránya.

\begin{center}\adjustbox{max width=\linewidth}{$\displaystyle \text{felső: }24+18=42\ \ldots\ 36+18=54\ \text{MHz}\qquad\frac{54}{42}=1{,}29$}\end{center}
\jegyzet{behelyettesítés}{}

\begin{center}\adjustbox{max width=\linewidth}{$\displaystyle \text{alsó: }24-18=6\ \ldots\ 36-18=18\ \text{MHz}\qquad\frac{18}{6}=3$}\end{center}
\jegyzet{behelyettesítés}{}

\begin{center}\adjustbox{max width=\linewidth}{$\displaystyle 1{,}29<3$}\end{center}
\jegyzet{ellenőrzés}{a felső keverés a jobb}

\lepes{b) Alsó keverés, 25,5 MHz}


\begin{center}\adjustbox{max width=\linewidth}{$\displaystyle f_\text{osc}=25{,}5-18=7{,}5\ \text{MHz}$}\end{center}
\jegyzet{behelyettesítés}{}

\lepes{c) A tükör}
Az a másik frekvencia, amely 7,5 MHz-cel keverve szintén 18 MHz-et ad. A különbség $f-7{,}5=\pm18$ ad 25,5-et (a hasznos jelet) és $-10{,}5$-et, ami a 10,5 MHz-es jelnek felel meg: ezt az \textbf{összegfrekvencia} viszi a KF-re.

\begin{center}\adjustbox{max width=\linewidth}{$\displaystyle f_t=|{\ca{f_v}}-2f_\mathrm{KF}|=|25{,}5-36|=10{,}5\ \text{MHz}$}\end{center}
\jegyzet{behelyettesítés}{}

\begin{center}\adjustbox{max width=\linewidth}{$\displaystyle 10{,}5+7{,}5=18\ \text{MHz}$}\end{center}
\jegyzet{ellenőrzés}{tényleg a KF-re esik}

\lepes{c) Mi vele a gond?}
A tükrök a sáv bejárásakor $|{\ca{24}}-36|=12$ MHz és $|36-36|=0$ között mozognak: ezt a keverő előtt egy hangolt szűrőnek kell elnyomnia. Ráadásul alsó keverésnél az oszcillátor a sáv tetején 18 MHz-en, pont a KF-en jár, és közvetlenül bejut a KF-erősítőbe. Felső keverésnél ($f_t=60\ldots72$ MHz) ezek a gondok nincsenek.
\begin{tipp}[A gyűjtemény megoldása]
A gyűjtemény tükörként $F+f_0=25{,}5+7{,}5=33$ MHz-et ad.
\end{tipp}
\begin{figyelem}[Javítás]
A 33 MHz a keverés összegfrekvenciája, nem tükör: $|33-7{,}5|=25{,}5\ne18$. A tükör 10,5 MHz.
\end{figyelem}
\begin{eredmeny}

\begin{center}\adjustbox{max width=\linewidth}{$\displaystyle \text{a) felső: }42\ldots54\ \text{MHz, alsó: }6\ldots18\ \text{MHz, a felső a jobb}\qquad\text{b) }7{,}5\ \text{MHz}\qquad\text{c) }f_t=10{,}5\ \text{MHz}$}\end{center}
\end{eredmeny}


% ------------------------------------------------------------------
\section{\texorpdfstring{8.3. feladat: DVB-T.HU multiplex nettó bitsebessége}{8.3. feladat: DVB-T.HU multiplex nettó bitsebessége}}\label{o8_03}
\setcounter{lepes}{0}
\begin{feladat}{8.3. feladat}
6817 vivő, ebből 6048 hasznos, $T_0=896$ \textmu{}s, $GI=1/4$, 64-QAM, RS(204,188), FEC 3/4. Mekkora a nettó bitsebesség, és hány 1,5 Mbit/s-os csatorna fér el?

{\footnotesize\sffamily Forrás: 5-ösöm 25\ldots{}26. oldal $\cdot$ 5o\_25\_2. Címkék: vizsga}
\end{feladat}
\begin{keresunk}
\begin{itemize}
\item $v_\text{nettó}$ (nettó bitsebesség) és a csatornák száma
\end{itemize}
\end{keresunk}
\begin{jeloles}
\begin{jeltabla}
$\ca{N}$ & az OFDM-vivők száma & $6817$\\
$N_\text{hasznos}$ & az adatot vivő vivők száma & \\
$\cb{T_0}$ & hasznos szimbólumidő & $896$\,\textmu{}s\\
$GI$ & a védőintervallum aránya & \\
$v_\text{nettó}$ & nettó bitsebesség & \emph{keresett}\\
$B$ & a csatorna sávszélessége & \\
$R_s$ & szimbólumsebesség vivőnként & \\
\end{jeltabla}
\par\smallskip\textbf{Adott még:}
\begin{itemize}
\item FEC 3/4 és RS(204,188): a két hibajavító kód aránya
\end{itemize}
\end{jeloles}

\lepes{Sávszélesség (ellenőrzésképpen)}
A vivőköz $1/{\cb{T_0}}=1{,}116$ kHz, a vivők sávja kb. $({\ca{N}}+1)/{\cb{T_0}}$ (A2).

\begin{center}\adjustbox{max width=\linewidth}{$\displaystyle B=\frac{6818}{896\ \mu\text{s}}=7{,}61\ \text{MHz}$}\end{center}
\jegyzet{behelyettesítés}{belefér a 8 MHz-es csatornába}

\lepes{Szimbólumsebesség}
Egy OFDM-szimbólum a védőintervallummal $(1+1/4)\cdot{\cb{896}}=1120$ \textmu{}s. Ennyi idő alatt 6048 hasznos QAM-szimbólum megy át.

\begin{center}\adjustbox{max width=\linewidth}{$\displaystyle R_s=\frac{6048}{1120\ \mu\text{s}}=5{,}4\ \text{MBd}$}\end{center}
\jegyzet{behelyettesítés}{}

\lepes{Nettó bitsebesség}
Szorozzuk a bit/szimbólummal, majd a két kód arányával.

\begin{center}\adjustbox{max width=\linewidth}{$\displaystyle 5{,}4\cdot6=32{,}4\ \text{Mbit/s}$}\end{center}
\jegyzet{számérték}{bruttó}

\begin{center}\adjustbox{max width=\linewidth}{$\displaystyle 32{,}4\cdot\frac34\cdot\frac{188}{204}=24{,}3\cdot0{,}922=22{,}4\ \text{Mbit/s}$}\end{center}
\jegyzet{számérték}{nettó}

\lepes{Csatornák}


\begin{center}\adjustbox{max width=\linewidth}{$\displaystyle \frac{22{,}4}{1{,}5}=14{,}9\ \Rightarrow\ 14\ \text{csatorna}$}\end{center}
\jegyzet{számérték}{egész számú csatorna fér el}
\begin{eredmeny}

\begin{center}\adjustbox{max width=\linewidth}{$\displaystyle v_\text{nettó}=22{,}4\ \text{Mbit/s}\qquad14\ \text{teljes csatorna (14,9)}$}\end{center}
\end{eredmeny}


% ------------------------------------------------------------------
\section{\texorpdfstring{8.4. feladat: Miért QAM és AM-VSB?}{8.4. feladat: Miért QAM és AM-VSB?}}\label{o8_04}
\setcounter{lepes}{0}
\begin{feladat}{8.4. feladat}
Miért használunk a műsorszórásban QAM-et és AM-VSB-t?

{\footnotesize\sffamily Forrás: 5-ösöm 24. oldal $\cdot$ 5o\_24\_1. Címkék: beugró}
\end{feladat}
\begin{keresunk}
\begin{itemize}
\item a két modulációs mód indoklása
\end{itemize}
\end{keresunk}
\begin{jeloles}
\par\smallskip\textbf{Adott még:}
\begin{itemize}
\item az analóg színes TV jelei: széles sávú, egyenszintet is tartalmazó képjel ($Y$) és két színkülönbségi jel
\end{itemize}
\end{jeloles}

\lepes{QAM}
A színsegédvivőt két, egymásra merőleges (koszinusz és szinusz) vivővel modulálják: két független jelet, a két színkülönbségi jelet visz ugyanabban a sávban, egymás zavarása nélkül. Ugyanez az elv a digitális QAM-nél: két kvadratúra-csatorna egy sávban.

\lepes{AM-VSB}
A képjel 5 MHz körüli széles és egyenszintet (0 Hz körüli összetevőt) is tartalmaz. AM-DSB-vel a sáv 10 MHz lenne: pazarló. SSB-hez a vivő közvetlen közelében kellene éles szűrő, ez DC-tartalmú jelnél nem valósítható meg. A csonka oldalsáv (VSB) az egyik oldalsávot nagyrészt levágja, csak egy kis maradékot hagy: a sáv majdnem a felére csökken, és a jel egyszerű burkolódetektorral demodulálható.
\begin{eredmeny}

\begin{center}\adjustbox{max width=\linewidth}{$\displaystyle \text{QAM: két független jel egy sávban}\qquad\text{AM-VSB: majdnem fele sáv, burkolódetektorral vehető}$}\end{center}
\end{eredmeny}
\begin{tipp}[Ellenőrzés]
\begin{itemize}
\item A gyűjtemény válasza („a QAM két szín esetén hasznos, frekvencianyereséggel jár (?)”) bizonytalan, a VSB-részlet nagyjából jó.
\end{itemize}
\end{tipp}
